\chapter{Introduction}
\label{ch:intro}

\section{Context}
\label{sec:context}

A growing share of what people ask a phone to do is not a question but an
\emph{action}: set an alarm, dim the screen, start a timer, turn on Do Not
Disturb, call a contact. These requests are short, frequent, and repetitive, and
they share a property that the conversational use of assistants does not. They
have a single correct outcome. Either the alarm is set for 7:00 or it is not.

How these actions are served is in the middle of an industry shift. The dominant
pattern routes a request to a large cloud model; the emerging one keeps as much
as possible \emph{on the device}. Both Google (Gemini Nano through Android
AICore) and Apple (Apple Intelligence, with Private Cloud Compute for overflow)
now run a small model locally and escalate only when they must. The motivation is
threefold: \emph{privacy}, because an action request often contains personal
context that should never leave the device; \emph{cost and latency}, because a
local model answers in milliseconds without a network round-trip or a per-call
bill; and \emph{availability}, because the device should still work offline. The
open engineering question this shift raises is the subject of this thesis: how
much language model does reliable on-device \emph{action} actually require?

\section{The problem}
\label{sec:problem}

The difficulty is a squeeze from both ends. A model small enough to run
comfortably on a phone is, on its own, an unreliable executor: asked to carry out
a device action by free-form reasoning, it either fails to act at all or invents
a command from its pretraining priors instead of the one the device expects
(Chapter~\ref{ch:eval} measures both failures). A model large enough to reason
reliably does not fit on the device, answers an order of magnitude more slowly,
costs money per call, and, as this thesis shows, is \emph{still} wrong on the
device-specific conventions that no amount of reasoning can recover.

So neither obvious option is satisfactory. Shipping a big model defeats the point
of going on-device; shipping a small free-form model is unreliable. The
interesting design space is between them: a tiny model that does not have to be
clever, paired with something that supplies the knowledge it lacks.

\section{The gap}
\label{sec:gap}

Two gaps in the existing picture motivate the work. First, the literature treats
``the model is too small'' as a single problem, when on a realistic device API it
is really two. One is \emph{discovery}: finding which setting, namespace, or key a
request maps to, which a larger model solves by exploring the API. The other is
\emph{convention}: knowing that ``total silence'' is encoded as the integer
\texttt{2}, or that ``largest font'' is the string \texttt{1.30}. Convention is
not discoverable from the API at all, because the meaning lives in a device's
accumulated conventions rather than in any queryable surface. These two
difficulties respond to completely different remedies, and conflating them hides
the most useful result in the thesis (Figure~\ref{fig:discovery}).

Second, the standard tool for compressing a capable model into a deployable one
is knowledge distillation into a smaller network's weights. That approach yields
an opaque student that must be retrained to extend. For an on-device system that
skill authors must extend and users must trust, a different distillation target
is attractive: compressing the teacher's knowledge into an \emph{inspectable,
plug-and-play symbolic artifact}. Whether such symbolic distillation is
sufficient to make a tiny model a reliable executor has not, to our knowledge,
been measured end-to-end.

\section{Thesis statement and research questions}
\label{sec:statement}

This thesis advances and defends a single claim:

\begin{quote}
On the repeatable, deterministic \emph{head} of on-device tasks, \textbf{distillation
substitutes for scale}: a 2-billion-parameter executor handed a distilled
symbolic pattern matches a $17\times$ larger free-form agent in binding accuracy,
at one on-device forward pass, because the bottleneck is not reasoning but device
knowledge the API does not expose. The remaining open limit is \emph{routing},
not execution.
\end{quote}

The claim is decomposed into five research questions, each answered in
Chapter~\ref{ch:eval}:

\begin{description}[leftmargin=2.2em,style=nextline]
  \item[RQ1 (scale vs.\ distillation)] How much model does reliable binding
    (intent $\rightarrow$ correct command) require, and can a distilled pattern
    close the gap a small model has?
  \item[RQ2 (the spectrum)] Does the value of the reflexer change with API
    difficulty: when is it an efficiency win and when an accuracy win?
  \item[RQ3 (the limit of scale)] Are there task difficulties that model scale
    cannot solve but a distilled pattern can?
  \item[RQ4 (teacher quality)] In teacher$\rightarrow$student distillation, how
    much does the choice of teacher (cloud vs.\ local) propagate to the on-device
    student?
  \item[RQ5 (coverage and routing)] What fraction of real on-device use-cases
    does the approach address, and where is the end-to-end bottleneck?
\end{description}

\section{Contributions}
\label{sec:contributions}

\begin{enumerate}[leftmargin=2em]
  \item \textbf{A two-tier local-first architecture} (A.L.F.R.E.D.) and the
    \emph{executioner paradigm}: a small reflexer treated as a constrained
    slot-filler, not a generator, with a heavier thinker reserved for the tail
    (Chapter~\ref{ch:design}).
  \item \textbf{Symbolic distillation}: knowledge distillation into a
    plug-and-play JSON pattern rather than student weights, with the
    teacher$\rightarrow$student loop measured end-to-end (Section~\ref{sec:p5}).
  \item \textbf{The discovery-vs-convention distinction}: two kinds of task
    difficulty, separated by what model scale can and cannot fix, demonstrated on
    a realistic device API (Section~\ref{sec:p4}).
  \item \textbf{Two end-state-oracle benchmarks} (SET-1 clean, SET-2
    Android-realistic) and a standardized binding-only harness, with an integrity
    protocol that forbids fabricated numbers and benchmark leakage into distilled
    artifacts (Chapter~\ref{ch:method}).
  \item \textbf{A coverage denominator} anchored to real Pixel and Apple
    on-device use-cases, and an honest localization of the open bottleneck at
    routing rather than execution (Sections~\ref{sec:coverage}--\ref{sec:routing}).
\end{enumerate}

\section{Thesis structure}
\label{sec:structure}

Chapter~\ref{ch:design} presents A.L.F.R.E.D.\ as a system: the two tiers, the
pattern artifact, the skill model, and the distillation step that connects them.
Chapter~\ref{ch:method} specifies how the system is evaluated, covering the
binding-only methodology, the end-state oracle, the benchmarks, and the integrity
rules. Chapter~\ref{ch:eval} is the core, establishing the five results above in
turn. Chapter~\ref{ch:discussion} draws out the implications for design and
productization and states the threats to validity. Chapter~\ref{ch:conclusion}
summarizes the contributions and outlines the next steps: device grounding,
deterministic value transforms, and the routing work that the evaluation
identifies as the remaining limit.
